\scanpage{165}\mspage{110}
(interprétés comme flèches dans $(\unsure{M|\mathcal M}_{/S})W_S^{-1}$) soient eux-mêmes fonctoriels en $X$.

On désigne par
\[
\mathrm{Ho}_W^{\mathrm{lc}}(S)\qquad\text{(lc pour ``loc. constant'').}
\]
la catégorie de fractions $(M_{/S})W_S^{-1}$, équivalente à $(\underline{\mathrm{Fib}}\,S)W_S^{-1}$.

\textcircled{II} Donnons nous un localiseur $W_S^{\unsure{0|o|\circ}}$ de $M_{/S}$
\[
W_S^0\subset W_S\subset\mathrm{Fl}(M_{/S}),
\]
avec \unsure{les hypothèses|la hypothèse} formelles similaires à celles de $W$ (cf 1.). On suppose que
\[
W_S^0\cap\mathrm{Fl}(\underline{\mathrm{Fib}}\,S)=W_S\mid\underline{\mathrm{Fib}}\,S.\tag*{3°)}
\]

(Cf pages 94--96 pour une façon générale d'obtenir de tels localiseurs. Il faut bien voir que pour définir des ``types ``d'homotopie'' relatifs sur $S$'', le localiseur $W_S$ est beaucoup trop grossier, \unsure{manifestement|moralement|évidemment} il \inserted{ne} donne qu'une catégorie des types d'hom. rel. localement constants, \unsure{\})} et c'est justement ça qui est montré par \deleted{\unsure{l'}} l'équivalence de \textcircled{I}, où $\underline{\mathrm{Fib}}\,S$ est visualisé comme formé des ``fibrations loc. triviales'' sur $S$.)
On a alors un deuxième localisé
\[
\mathrm{Ho}_{W_S^0}(S)\overset{\mathrm{def}}{=}(M_{/S})(W_S^0)^{-1},
\]
et deux foncteurs
\diagram{\begin{tikzcd}[ampersand replacement=\&]
{(*)} \& \mathrm{Ho}_{W_S^0}(S) \arrow[r,shift left,"\rho"] \& \mathrm{Ho}_W^{\mathrm{lc}}(S) \arrow[l,shift left,"i"]
\end{tikzcd}}
où $\rho$ est le foncteur de localisation \inserted{(strict)} évident, dû à l'inclusion $W_S^0\subset W_S$, et où $i$ est
% archive sheet number: 164
